\documentclass[a4paper,12pt]{article}

\usepackage[T2A]{fontenc}
\usepackage[utf8]{inputenc}
\usepackage[russian]{babel}

\usepackage{amsmath,amssymb,mathtools}
\usepackage{helvet}
\renewcommand{\familydefault}{\sfdefault}

\usepackage{icomma}
\usepackage[a4paper,left=22mm,right=22mm,top=20mm,bottom=22mm]{geometry}
\usepackage{microtype}
\usepackage{tikz}
\usetikzlibrary{calc}
\usepackage{tkz-euclide}
\usepackage{pgfplots}
\pgfplotsset{compat=1.18}
\usepackage{tabularx,booktabs}
\usepackage{enumitem}
\usepackage{hyperref}
\usepackage{parskip}
\usepackage{fancyhdr}
\usepackage{setspace}
\usepackage{xcolor}

\definecolor{accent}{HTML}{008080}
\definecolor{darkaccent}{HTML}{004D4D}
\definecolor{textgray}{HTML}{333333}
\definecolor{softgray}{HTML}{F2F4F4}

\hypersetup{
    colorlinks=true,
    linkcolor=darkaccent,
    urlcolor=darkaccent
}

\setlength{\parindent}{0pt}
\setlength{\headheight}{15pt}
\onehalfspacing

\pagestyle{fancy}
\fancyhf{}
\fancyhead[L]{\small\color{textgray}Повторение алгебры и производная}
\fancyhead[R]{\small\color{textgray}Николь Саркисьянц}
\fancyfoot[C]{\small\thepage}
\renewcommand{\headrulewidth}{0.3pt}
\renewcommand{\headrule}{\hbox to\headwidth{\color{accent}\leaders\hrule height \headrulewidth\hfill}}

\newcommand{\lessonsection}[1]{%
    \vspace{0.5em}
    {\Large\bfseries\color{darkaccent}#1}\par
    \vspace{0.2em}
    {\color{accent}\rule{\linewidth}{0.6pt}}
    \vspace{0.5em}
}

\newcommand{\important}[1]{\textbf{\color{darkaccent}#1}}

\begin{document}

% =========================================================
% ТИТУЛЬНЫЙ ЛИСТ
% =========================================================

\begin{titlepage}
\thispagestyle{empty}

\begin{center}

\vspace*{2.4cm}

{\Large\color{accent}\textbf{ЧЕК-ЛИСТ}}

\vspace{0.7cm}

{\Huge\bfseries
Повторение алгебры\\[0.2cm]
и введение в производную}

\vspace{1.2cm}

{\large
Движение по воде, степени и корни,\\
функциональные подстановки, алгебраические дроби,\\
базовые правила дифференцирования}

\vfill

{\large
Лёвин Артём Александрович\\[0.15cm]
\textbf{эксклюзивно для Николь Саркисьянц}
}

\vspace{0.5cm}

{\large \today}

\vspace{1.6cm}

\begin{tikzpicture}[x=1cm,y=1cm]
    \draw[
        accent,
        line width=1.1pt
    ]
    (0,0)
    .. controls (3,0.75) and (7,-0.75) ..
    (10,0)
    .. controls (12,0.35) and (14,0.35) ..
    (16,0);
\end{tikzpicture}

\vspace{0.8cm}

\end{center}
\end{titlepage}

% =========================================================
% ОСНОВНОЙ ЧЕК-ЛИСТ
% =========================================================

\lessonsection{1. Задачи на движение по воде}

Для задач на движение удобно сразу строить таблицу
\[
\text{скорость}\qquad \text{время}\qquad \text{расстояние}.
\]

Если собственная скорость теплохода равна \(v\), а скорость течения равна \(u\), то

\[
v_{\text{по течению}}=v+u,
\qquad
v_{\text{против течения}}=v-u.
\]

Следовательно,
\[
t_{\text{по течению}}=\frac{s}{v+u},
\qquad
t_{\text{против течения}}=\frac{s}{v-u}.
\]

Обязательно:
\[
v>u,
\]
иначе движение против течения физически невозможно.

\important{Если в условии есть стоянка}, её время также входит в общее время путешествия.

Например, если теплоход проходит по \(200\) км в каждую сторону, скорость течения равна \(15\) км/ч, стоянка занимает \(10\) часов, а всё путешествие --- \(40\) часов, то

\[
\frac{200}{v+15}+10+\frac{200}{v-15}=40,
\qquad v>15.
\]

После переноса времени стоянки:

\[
\frac{200}{v+15}+\frac{200}{v-15}=30.
\]

Общий знаменатель:
\[
(v+15)(v-15).
\]

Здесь обязательно используется формула разности квадратов:

\[
(a+b)(a-b)=a^2-b^2.
\]

Поэтому
\[
(v+15)(v-15)=v^2-225.
\]

Числитель:
\[
200(v-15)+200(v+15)=400v.
\]

Получаем
\[
\frac{400v}{v^2-225}=30.
\]

После упрощения:
\[
3v^2-40v-675=0.
\]

Дискриминант:
\[
D=(-40)^2-4\cdot3\cdot(-675)
=1600+8100
=9700.
\]

\[
\sqrt{9700}=10\sqrt{97}.
\]

Корни:
\[
v_{1,2}=\frac{40\pm10\sqrt{97}}{6}.
\]

С учётом условия \(v>15\):
\[
v=\frac{20+5\sqrt{97}}{3}.
\]

\important{Типичная ошибка:} раскрывать
\[
(v+15)(v-15)
\]
без распознавания формулы
\[
(a+b)(a-b)=a^2-b^2.
\]

\vspace{0.4em}

\textbf{Формулы сокращённого умножения, которые необходимо помнить:}
\[
(a+b)^2=a^2+2ab+b^2,
\]
\[
(a-b)^2=a^2-2ab+b^2,
\]
\[
(a-b)(a+b)=a^2-b^2,
\]
\[
a^3-b^3=(a-b)(a^2+ab+b^2),
\]
\[
a^3+b^3=(a+b)(a^2-ab+b^2).
\]

% =========================================================

\lessonsection{2. Десятичные дроби и степени числа \(10\)}

Десятичную дробь часто удобно представить в виде произведения целого числа и степени \(10\).

Например,
\[
1{,}23=123\cdot10^{-2},
\]
\[
45{,}7=457\cdot10^{-1},
\]
\[
0{,}457=457\cdot10^{-3}.
\]

\important{Принцип:} если запятую переносим вправо на \(k\) разрядов, появляется множитель
\[
10^{-k}.
\]

Это особенно полезно в дробных выражениях, поскольку одинаковые числовые множители можно сократить, а степени числа \(10\) объединить по правилам:

\[
10^a\cdot10^b=10^{a+b},
\]
\[
\frac{10^a}{10^b}=10^{a-b}.
\]

Например,
\[
\frac{
(123\cdot10^{-2})(457\cdot10^{-1})
}{
(123\cdot10^{-1})(457\cdot10^{-3})
}.
\]

Сокращаем \(123\) и \(457\):

\[
10^{-2-1-(-1)-(-3)}
=
10^1
=
10.
\]

\important{Типичная ошибка:} неверно работать с двойным минусом:
\[
-2-(-3)=-2+3=1.
\]

% =========================================================

\lessonsection{3. Функция и подстановка выражения вместо аргумента}

Если задана функция
\[
P(x)=x^2-x-6,
\]
то для вычисления \(P(2)\) нужно \textbf{везде вместо \(x\)} подставить число \(2\):

\[
P(2)=2^2-2-6=-4.
\]

Если вместо аргумента подставляется выражение, принцип остаётся тем же.

Например,
\[
P(2x-3)
=
(2x-3)^2-(2x-3)-6.
\]

\important{Главная идея:}
\[
P(\boxed{\phantom{x}})
\]
означает, что содержимое скобок подставляется на место аргумента во всей формуле функции.

Рассмотрим другой пример:
\[
P(b)=
\left(b+\frac{3}{b}\right)
\left(3b+\frac{1}{b}\right),
\qquad b\ne0.
\]

Тогда
\[
P\left(\frac1b\right)
=
\left(
\frac1b+
\frac{3}{1/b}
\right)
\left(
3\cdot\frac1b+
\frac{1}{1/b}
\right).
\]

Используем
\[
\frac{3}{1/b}=3b,
\qquad
\frac{1}{1/b}=b.
\]

Получаем
\[
P\left(\frac1b\right)
=
\left(\frac1b+3b\right)
\left(\frac3b+b\right).
\]

Причём
\[
b+\frac3b=\frac{b^2+3}{b}
=
\frac3b+b,
\]
а также
\[
3b+\frac1b=\frac{3b^2+1}{b}
=
\frac1b+3b.
\]

Поэтому
\[
P(b)=P\left(\frac1b\right)
\]
и
\[
\frac{P(b)}{P(1/b)}=1
\]
при условии, что знаменатель этой дроби определён и отличен от нуля.

\important{ОДЗ:}
\[
b\ne0,
\]
а при делении дополнительно требуется
\[
P(1/b)\ne0.
\]

% =========================================================

\lessonsection{4. Алгебраические дроби: основной метод}

Если дано выражение
\[
\frac{5\sqrt{x}+2}{\sqrt{x}}
-
\frac{2\sqrt{x}}{x},
\]
то сначала фиксируем ОДЗ.

Из знаменателей:
\[
\sqrt{x}\ne0,
\qquad x\ne0.
\]

Кроме того, корень должен существовать:
\[
x\ge0.
\]

Итого:
\[
x>0.
\]

\important{Основной метод: привести дроби к общему знаменателю.}

Общий знаменатель:
\[
x\sqrt{x}.
\]

Тогда
\[
\frac{5\sqrt{x}+2}{\sqrt{x}}
-
\frac{2\sqrt{x}}{x}
=
\frac{x(5\sqrt{x}+2)-2x}{x\sqrt{x}}.
\]

Раскрываем скобки:
\[
=
\frac{5x\sqrt{x}+2x-2x}{x\sqrt{x}}.
\]

Получаем
\[
=
\frac{5x\sqrt{x}}{x\sqrt{x}}
=5.
\]

Сокращение допустимо, поскольку на ОДЗ
\[
x\sqrt{x}\ne0.
\]

\important{Запрещённый приём:} нельзя произвольно возводить отдельные части выражения в квадрат. Возведение обеих частей в квадрат является преобразованием \textbf{уравнения}, а не способом упрощения произвольной суммы или разности.

% =========================================================

\lessonsection{5. Почленное деление}

Дополнительный способ упрощения дроби:

\[
\frac{A+B}{C}
=
\frac AC+\frac BC,
\qquad C\ne0.
\]

Аналогично:
\[
\frac{A-B}{C}
=
\frac AC-\frac BC.
\]

Например,
\[
\frac{5\sqrt{x}+2}{\sqrt{x}}
=
5+\frac{2}{\sqrt{x}}.
\]

Во второй дроби используем
\[
x=\sqrt{x}\cdot\sqrt{x},
\qquad x\ge0.
\]

При \(x>0\):
\[
\frac{2\sqrt{x}}{x}
=
\frac{2\sqrt{x}}
{\sqrt{x}\sqrt{x}}
=
\frac2{\sqrt{x}}.
\]

Следовательно,
\[
5+\frac2{\sqrt{x}}-\frac2{\sqrt{x}}=5.
\]

\important{Приоритетный универсальный метод} для нескольких алгебраических дробей --- приведение к общему знаменателю. Почленное деление удобно использовать как дополнительный приём.

% =========================================================

\lessonsection{6. Корни, дроби и единый степенной вид}

При работе со степенями, корнями и дробями полезно переводить выражение в единый степенной вид.

\textbf{1. Составные числа}

Например,
\[
32=2^5,
\qquad
81=3^4.
\]

\textbf{2. Корень как степень}

\[
\sqrt[k]{a^n}=a^{n/k}
\]
на области, где выражение корректно определено.

В частности,
\[
\sqrt[k]{a}=a^{1/k}.
\]

Например,
\[
\sqrt[5]{\sqrt[3]{a}}
=
\left(a^{1/3}\right)^{1/5}
=
a^{1/15}.
\]

\textbf{3. Дробь как отрицательная степень}

\[
\frac1{a^n}=a^{-n},
\qquad a\ne0.
\]

Таким образом, сложные выражения часто становятся набором обычных степеней.

Например,
\[
\frac{
15\sqrt[5]{\sqrt[3]{a}}
-
7\sqrt[5]{\sqrt[3]{a}}
}{
2\sqrt[5]{\sqrt[3]{a}}
}.
\]

Обозначим
\[
t=\sqrt[5]{\sqrt[3]{a}},
\qquad t\ne0.
\]

Тогда
\[
\frac{15t-7t}{2t}
=
\frac{8t}{2t}
=
4.
\]

Можно также предварительно записать
\[
t=a^{1/15}.
\]

\important{Полезная стратегия:}
если один и тот же сложный фрагмент многократно повторяется, рассмотрите замену
\[
t=\text{повторяющееся выражение}.
\]

% =========================================================

\lessonsection{7. Производная: обозначение и первая формула}

Производная функции \(f(x)\) обозначается
\[
f'(x).
\]

Для степенной функции действует правило
\[
\boxed{
(x^n)'=nx^{n-1}
}
\]
для тех \(x\), где соответствующая функция и её производная определены.

Алгоритм:

\begin{enumerate}[label=\arabic*)]
    \item показатель степени переносим в множитель;
    \item показатель степени уменьшаем на \(1\).
\end{enumerate}

Например,
\[
f(x)=x^2
\quad\Longrightarrow\quad
f'(x)=2x^{2-1}=2x.
\]

\[
f(x)=x^3
\quad\Longrightarrow\quad
f'(x)=3x^2.
\]

\[
f(x)=x^4
\quad\Longrightarrow\quad
f'(x)=4x^3.
\]

% =========================================================

\lessonsection{8. Отрицательные и дробные показатели}

Правило степенной функции распространяется и на многие отрицательные и дробные показатели.

Например,
\[
f(x)=x^{-2}.
\]

Тогда
\[
f'(x)
=
-2x^{-3}
=
-\frac2{x^3},
\qquad x\ne0.
\]

Если
\[
f(x)=x^{1/2}=\sqrt{x},
\]
то
\[
f'(x)
=
\frac12x^{-1/2}
=
\frac{1}{2\sqrt{x}},
\qquad x>0.
\]

Обратите внимание: исходная функция \(\sqrt{x}\) определена при
\[
x\ge0,
\]
но формула её производной содержит \(\sqrt{x}\) в знаменателе, поэтому производная в данном виде определена при
\[
x>0.
\]

\important{Перед дифференцированием часто полезно выполнить преобразование:}
\[
\frac1{x^m}=x^{-m},
\]
\[
\sqrt[k]{x^m}=x^{m/k}.
\]

% =========================================================

\lessonsection{9. Производная постоянной}

Если функция равна постоянному числу,
\[
f(x)=C,
\]
то
\[
\boxed{f'(x)=0}.
\]

Например,
\[
(5)'=0,
\]
\[
(-100)'=0,
\]
\[
(\pi)'=0,
\]
\[
(\sqrt3)'=0.
\]

Причина проста: эти величины не зависят от \(x\).

Важно различать:
\[
\sqrt3
\quad\text{и}\quad
\sqrt{x}.
\]

Первое выражение является числом, второе зависит от \(x\).

% =========================================================

\lessonsection{10. Постоянный множитель}

Если постоянное число умножается на функцию, его можно оставить перед производной:

\[
\boxed{
(Cf(x))'=C f'(x)
}
\]

Например,
\[
f(x)=5x^2.
\]

Тогда
\[
f'(x)
=
5(x^2)'
=
5\cdot2x
=
10x.
\]

Ещё пример:
\[
f(x)=\sqrt5\,x^3.
\]

Число \(\sqrt5\) является постоянным множителем:

\[
f'(x)
=
\sqrt5\,(x^3)'
=
3\sqrt5\,x^2.
\]

Если
\[
f(x)=\pi x^{1/4},
\]
то
\[
f'(x)
=
\pi\cdot\frac14x^{-3/4}
=
\frac{\pi}{4x^{3/4}},
\qquad x>0.
\]

% =========================================================

\lessonsection{11. Производная суммы и разности}

Для суммы:
\[
\boxed{
(f(x)+g(x))'
=
f'(x)+g'(x)
}
\]

Для разности:
\[
\boxed{
(f(x)-g(x))'
=
f'(x)-g'(x)
}
\]

То есть производную считаем от каждого слагаемого отдельно.

Например,
\[
f(x)=5+x^2.
\]

Тогда
\[
f'(x)
=
(5)'+(x^2)'
=
0+2x
=
2x.
\]

Ещё пример:
\[
f(x)=7-3x^4+\sqrt{x}.
\]

Тогда
\[
f'(x)
=
0-12x^3+\frac1{2\sqrt{x}},
\qquad x>0.
\]

% =========================================================

\lessonsection{12. Когда степенное правило пока неприменимо напрямую}

Формула
\[
(x^n)'=nx^{n-1}
\]
в рассматриваемом виде требует, чтобы показатель \(n\) был постоянным числом.

Например,
\[
x^5,\qquad
x^{-2},\qquad
x^{1/3}
\]
подходят непосредственно.

Выражение
\[
x^{2x+3}
\]
имеет переменный показатель степени.

К нему простое правило
\[
(x^n)'=nx^{n-1}
\]
применять нельзя.

Для таких функций используются дополнительные методы дифференцирования.

% =========================================================

\lessonsection{13. Короткий алгоритм работы с производной}

Перед вычислением производной:

\begin{enumerate}[label=\arabic*)]
    \item Упростите вид функции, если это возможно.
    \item Переведите дробные и корневые степени в степенной вид.
    \item Отделите постоянные множители.
    \item При сложении и вычитании работайте с каждым слагаемым отдельно.
    \item Примените
    \[
    (x^n)'=nx^{n-1}.
    \]
    \item Производную постоянной замените на \(0\).
    \item Проверьте ОДЗ исходной функции и полученной производной.
\end{enumerate}

% =========================================================

\lessonsection{14. Типичные ошибки}

\begin{enumerate}[label=\arabic*)]
    \item Забывать формулу
    \[
    (a-b)(a+b)=a^2-b^2.
    \]

    \item Сокращать слагаемые вместо множителей.

    Например, в
    \[
    \frac{x+3}{x}
    \]
    нельзя ``сократить \(x\)''.

    \item Игнорировать ОДЗ дробей и корней.

    \item При делении на дробь забывать переход к умножению на обратную:
    \[
    \frac{a}{1/b}=ab.
    \]

    \item Неверно раскрывать двойной минус:
    \[
    a-(-b)=a+b.
    \]

    \item Произвольно возводить части выражения в квадрат.

    \item В производной забывать уменьшить показатель степени на \(1\).

    \item Считать, что производная постоянного числа равна самому числу:
    \[
    C'\ne C,
    \qquad
    C'=0.
    \]

    \item При наличии постоянного множителя ошибочно дифференцировать его как переменную.

    \item Применять степенную формулу напрямую к
    \[
    x^{g(x)},
    \]
    где показатель степени зависит от \(x\).
\end{enumerate}

% =========================================================
% ТРЕНИРОВОЧНЫЙ БЛОК
% =========================================================

\clearpage
\lessonsection{Тренировочный блок}

\textbf{Уровень: 3 средних, 6 выше среднего, 1 сложное.}

\begin{enumerate}[label=\textbf{\arabic*.},leftmargin=*]

\item
Упростите:
\[
(7-a)(7+a).
\]

\item
Представьте без десятичных запятых с помощью степеней числа \(10\) и вычислите:
\[
\frac{1{,}28\cdot0{,}35}{12{,}8\cdot0{,}0035}.
\]

\item
Найдите производную:
\[
f(x)=7x^5-4.
\]

\item
Дана функция
\[
P(t)=
\left(t+\frac4t\right)
\left(2t+\frac1t\right),
\qquad t\ne0.
\]
Запишите в упрощённом виде
\[
P\left(\frac1t\right).
\]

\item
Упростите выражение и укажите ОДЗ:
\[
\frac{4\sqrt{x}+3}{\sqrt{x}}
-
\frac{3\sqrt{x}}{x}.
\]

\item
Упростите:
\[
\frac{
11\sqrt[4]{\sqrt[3]{a}}
-
5\sqrt[4]{\sqrt[3]{a}}
}{
3\sqrt[4]{\sqrt[3]{a}}
}.
\]
Укажите ограничения, необходимые для выполненного сокращения.

\item
Найдите производную функции:
\[
f(x)=
\frac5{x^3}
-
2\sqrt{x}
+
9.
\]

\item
Найдите производную:
\[
f(x)=
3\sqrt7\,x^{5/2}
-
\frac4{x^2}
+
\pi.
\]

\item
Теплоход проходит \(120\) км по течению и \(120\) км против течения. Скорость течения равна \(5\) км/ч. Всё движение заняло \(10\) часов.

Найдите собственную скорость теплохода.

\item
Упростите выражение
\[
\frac{
\displaystyle
\frac{6\sqrt{x}+5}{\sqrt{x}}
-
\frac{5\sqrt{x}}{x}
}{
\displaystyle
\frac{2}{x^{-1/2}}
}
\]
и затем найдите производную полученной функции. Обязательно укажите ОДЗ исходного выражения.

\end{enumerate}

% =========================================================
% ОТВЕТЫ
% =========================================================

\clearpage
\lessonsection{Ответы к тренировочному блоку}

\renewcommand{\arraystretch}{1.45}

\begin{table}[ht]
\centering
\caption{Краткие ответы}
\begin{tabularx}{\linewidth}{@{}cX@{}}
\toprule
\textbf{№} & \textbf{Ответ} \\
\midrule

1 &
\[
49-a^2.
\]
\\

2 &
\[
10.
\]
\\

3 &
\[
f'(x)=35x^4.
\]
\\

4 &
\[
P\left(\frac1t\right)
=
\left(\frac1t+4t\right)
\left(\frac2t+t\right),
\qquad t\ne0.
\]
\\

5 &
\[
4,
\qquad x>0.
\]
\\

6 &
\[
2.
\]
Для сокращения требуется
\[
\sqrt[4]{\sqrt[3]{a}}\ne0,
\]
то есть
\[
a\ne0.
\]
\\

7 &
Сначала
\[
f(x)=5x^{-3}-2x^{1/2}+9.
\]
Поэтому
\[
f'(x)
=
-15x^{-4}-x^{-1/2}
=
-\frac{15}{x^4}-\frac1{\sqrt{x}},
\qquad x>0.
\]
\\

8 &
\[
f'(x)
=
\frac{15\sqrt7}{2}x^{3/2}
+
\frac8{x^3}.
\]
Для исходной функции:
\[
x>0.
\]
\\

9 &
Пусть собственная скорость равна \(v\), \(v>5\):
\[
\frac{120}{v+5}+\frac{120}{v-5}=10.
\]
После преобразований:
\[
v^2-24v-25=0.
\]
\[
v=25\ \text{км/ч}
\]
(корень \(v=-1\) не соответствует смыслу задачи).
\\

10 &
ОДЗ исходного выражения:
\[
x>0.
\]

Числитель:
\[
\frac{6\sqrt{x}+5}{\sqrt{x}}
-
\frac{5\sqrt{x}}x
=
6.
\]

Знаменатель:
\[
\frac2{x^{-1/2}}
=
2\sqrt{x}.
\]

Поэтому
\[
f(x)=\frac6{2\sqrt{x}}
=
3x^{-1/2}.
\]

Производная:
\[
f'(x)
=
-\frac32x^{-3/2}
=
-\frac{3}{2x\sqrt{x}}.
\]
\\

\bottomrule
\end{tabularx}
\end{table}

\end{document}